% from ega4-2-fr.tex
\oldpage[IV-2]{198}

\item[\emph{ii)}] Nous avons déjà signalé (6.4.3) que lorsque $\boldsymbol P$ est
l'une des propriétés (ii) ou $(\mathrm{iii}_1)$ de (7.3.8), nous ignorons
si tout anneau local noethérien est un $\boldsymbol P$-anneau; rappelons
toutefois que lorsque $A$ est un \emph{anneau quotient d'un anneau de
Cohen--Macaulay}, les fibres formelles de $A$ sont des \emph{schémas de
Cohen--Macaulay} (6.3.8).
\item[\emph{iii)}] Le cas le plus intéressant de la notion de $\boldsymbol P$-anneau
est celui qui correspond à la propriété la plus forte (iv$'$) de
(7.3.8), c'est-à-dire les anneaux dont \emph{les fibres formelles sont
géométriquement régulières}. Les corps, et plus généralement les anneaux
locaux noethériens \emph{complets}, vérifient trivialement cette
propriété.
\item[\emph{iv)}] Soit $A$ un anneau local noethérien de \emph{dimension 1};
$\operatorname{Spec}(A)$ est alors formé du point fermé $a$,
correspondant à l'idéal maximal $\mathfrak m$, et des points maximaux
$b_i$ $(1\leqslant i\leqslant r)$ correspondant aux idéaux premiers
minimaux de $A$. On a $\dim(\widehat A)=1$ et l'idéal maximal de
$\widehat A$ est $\mathfrak m\widehat A$; la fibre formelle de $A$ au
point $a$ est donc $\operatorname{Spec}(k)$, où $k=A/\mathfrak m$ est le
corps résiduel de $A$; la fibre formelle en chacun des $b_i$ est le
spectre d'un anneau artinien dont les corps résiduels sont les corps
résiduels $L_{ij}$ de $\operatorname{Spec}(\widehat A)$ aux points
maximaux $b_{ij}$ au-dessus de $b_i$. Comme un anneau artinien est un
anneau de Cohen--Macaulay, on voit que $A$ est un
$\boldsymbol P$-anneau lorsque $\boldsymbol P$ est la propriété (ii) de
(7.3.8). En outre, comme un anneau artinien réduit est une somme directe
de corps, les propriétés suivantes sont équivalentes (6.7.7) :
\begin{enumerate}
\item[\emph{a)}] les fibres formelles de $A$ sont géométriquement réduites;
\item[\emph{b)}] les fibres formelles de $A$ sont géométriquement normales;
\item[\emph{c)}] les fibres formelles de $A$ sont géométriquement régulières;
\end{enumerate}
en outre, lorsque $A$ est \emph{réduit}, elles sont aussi équivalentes à :
\begin{enumerate}
\item[\emph{d)}] $\widehat A$ est réduit et $L_{ij}$ est une extension séparable de
$K_i$ pour tout couple $(i,j)$ (4.6.1).
\end{enumerate}
En particulier, si $A$ est un \emph{anneau de valuation discrète}, $K$
son corps des fractions, et si $\widehat K$ est le complété de $K$ pour
une valuation correspondant à $A$ (corps des fractions de $\widehat A$),
pour que les fibres formelles de $A$ soient géométriquement régulières,
il faut et il suffit que $\widehat K$ soit une \emph{extension séparable}
de $K$. Ce sera toujours le cas lorsque $K$ est de \emph{caractéristique
0}.
\end{enumerate}
\end{remarks}

\subsection{Permanence des propriétés des fibres formelles}
\label{subsection:IV.7.4-fr}

\begin{env}[7.4.1]
\label{IV.2.s7.4.1-fr}
Dans toute la suite du numéro, nous supposons que la propriété
$\boldsymbol P$ est de la forme définie dans (7.3.1), et \emph{vérifie
les conditions $(\mathrm P_{\mathrm I})$, $(\mathrm P_{\mathrm{II}})$ et
$(\mathrm P_{\mathrm{III}})$} de (7.3.4). En outre, nous supposons que la
propriété $\boldsymbol Q$ qui sert à définir $\boldsymbol P$ est telle
que, pour toute générisation $z'$ de $z$ dans $Z$,
$\boldsymbol Q(\mathscr O_z,k)$ entraîne
$\boldsymbol Q(\mathscr O_{z'},k)$.
\end{env}

\begin{lemma}[7.4.2]
\label{IV.2.s7.4.2-fr}
Soient $A$, $A'$ des anneaux locaux noethériens,
$\varphi:A\to A'$ un homomorphisme local tel que
$f={}^a\varphi:\operatorname{Spec}(A')\to\operatorname{Spec}(A)$ soit un
$\boldsymbol P$-morphisme. Alors, si les fibres formelles de $A'$ sont
géométriquement régulières, $A$ est un $\boldsymbol P$-anneau.
\end{lemma}

